\section{Related Work}

\paragraph{Agile Locomotion}
RL has demonstrated tremendous success in obtaining robust and adaptive controllers for legged robot~\citep{peng2020learning, fu2021minimizing,  kumar2021rma, rudin2022learning, aractingi2023controlling}.
This includes agile skills such as high-speed running~\citep{bellegarda2022robust,caluwaerts2023barkour, margolis2024rapid, he2024agile}, recovering from falling~\citep{hwangbo2019learning,yang2020multi,smith2022legged, inoue2024body}, jumping~\citep{bellegarda2020robust, margolis2021learning, li2023robust, smith2023learning, li2023learning, yang2023generalized, li2024reinforcement, zhang2024learning},
climbing obstacles~\cite{lee2020learning, siekmann2021blind,miki2022learning,hoeller2022neural,rudin2022advanced,agarwal2023legged,yang2023neural,loquercio2023learning, chane2024cat, cheng2024quadruped, hoeller2024anymal,zhuang2023robot,cheng2023parkour}
bipedal walking with quadruped robots~\citep{fuchioka2023opt,li2023learning,smith2023learning,cheng2023parkour,huang2024diffuseloco} and walking inside confined spaces~\citep{miki2024learning, xu2024dexterous, cheng2024quadruped}.
Learning multi-skill locomotion policies, as required in parkour, can be done by training separate policies for each skill, then coordinating them with a high-level planner~\citep{yang2020multi, kim2022humanconquad, huang2023creating, hoeller2024anymal} or distilling them into a single policy~\citep{caluwaerts2023barkour, zhuang2023robot, huang2024diffuseloco}.
Instead, we follow~\citep{cheng2023parkour} and learn multi-task policies directly through RL.

\paragraph{Safe Locomotion}
Safety mechanisms have been implemented to ensure safer outcome while performing agile skills~\citep{yang2022safe, he2024agile}.
Following~\citep{lee2023evaluation, kim2024not,chane2024cat}, we employ constraints alongside rewards in RL to deter undesired behaviors.
This not only allows for aggressive optimization of agility while ensuring safety but also simplifies the process of reward tuning for RL.
In particular, we exploit the reformulation of Constraints as Terminations~\citep{chane2024cat} (CaT), which was demonstrated as an overhead on top of on-policy RL algorithms~\citep{schulman2017proximal}, and which we extend to the off-policy formulation, more suitable to our case as explained below.


\paragraph{Sample-Efficient Learning}
\citep{bogdanovic2022model, smith2023learning} accelerate RL with demonstrations obtained from trajectory optimization.
Other works aimed at exploring RL methods sample efficient enough to train locomotion policies directly in the real-world~\cite{smith2022legged, wu2023daydreamer}.
Our approach repurposes many of these designs to bypass the computational cost of RL from pixels while obtaining a pure end-to-end RL method, unlocking several advantages exhibited below.

\paragraph{Vision-Based Locomotion}
Prior methods often separate perception from control using intermediate representations such as elevation maps~\citep{kleiner2007real,peng2017deeploco,fankhauser2018probabilistic,gangapurwala2021real,tsounis2020deepgait,rudin2022learning,miki2022learning, duan2023learning, hoeller2024anymal, chane2024cat}, traceability maps ~\citep{chavez2018learning, yang2021real} or visual odometry~\citep{bloesch2015robust,wisth2019robust,wisth2022vilens,buchanan2022learning},  for downstream planning and control~\citep{wermelinger2016navigation,mastalli2017trajectory,magana2019fast,jenelten2020perceptive,kim2020vision,agrawal2022vision}.
Recently, locomotion from pixels~\citep{yu2021visual,margolis2021learning,loquercio2023learning,agarwal2023legged,zhuang2023robot,cheng2023parkour} has emerged as a powerful paradigm that more tightly coordinates vision and control, often relying on teacher-student approaches (typically by distilling an observation-privileged policy), yet raising some limited action-perception behaviors.

